% This is a template for students in MATH 3000 at FSU to use for the final draft of their midterm exam.

% All of this stuff with '%' in front is a comment and ignored by the compiler.
%
% The lines before the "\begin{document}" line is called the preamble.
% This is where you load particular packages you need.
% Until you are more experienced, or the program says you are missing packages, it is safe to ignore most of the preamble.
%
%----------------------------------

\documentclass[12pt]{article}
\usepackage[margin=1in]{geometry}% Change the margins here if you wish.
\setlength{\parindent}{0pt} % This is the set the indent length for new paragraphs, change if you want.
\setlength{\parskip}{5pt} % This sets the distance between paragraphs, which will be used anytime you have a blank line in your LaTeX code.
\pagenumbering{gobble}% This means the page will not be numbered. You can comment it out if you like page numbers.

%------------------------------------

% These packages allow the most of the common "mathly things"
\usepackage{amsmath,amsthm,amssymb}

% This package allows you to add images.
\usepackage{graphicx}
\usepackage{float}

% These are theorem environments.  This should cover everything you need, and you should be able to tell what environment goes with what type of result, but please let me know if I've missed anything.
\newtheorem{euclidtheorem}{Proposition}[section]
\newtheorem{classtheorem}{Theorem}
\newtheorem{theorem}{Theorem}[section]
\newtheorem*{theorem*}{Theorem}


%These help to format the names of the results the way we are in class and in notes.
\renewcommand*{\theeuclidtheorem}{\Roman{section}.\arabic{theorem}}
\renewcommand*{\thetheorem}{\arabic{section}.\arabic{theorem}}
\renewcommand*{\theclasstheorem}{\Alph{theorem}}



% Should you need any additional packages, you can load them here. If you've looked up something (like on DeTeXify), it should specify if you need a special package.  Just copy and paste what is below, and put the package name in the { }.  
\usepackage{wasysym} %this lets me make smiley faces :-)

\title{Geometry Midterm}

% You are the author, put your name here.
\author{Sarah Wright}

% You can change the date to be something other than the current date if you want.
\date{\today}

\begin{document}

\maketitle

Write a brief summary of how your oral exam went.  Include some aspect that you worked hard on and went well.  Also describe something you will focus your practice on for the remainder of the semester and the final exam.  What you include below for each problem should be a formal write-up of final draft quality; follow the same Writing Style Guide as we use for papers.

% 
\begin{enumerate}
% PROBLEM #1 - CONSTRUCTION
\item  For my construction, I chose to construct a perpendicular.
\setcounter{section}{1}
\setcounter{theorem}{12}
% Change the introduction and counters appropriately to reflect which construction you chose.  The intro should be in modern language, and you can copy the language of the proposition from Euclid.
\begin{euclidtheorem}
To a given infinite straight line, from a given point which is not on it, to draw a perpendicular straight line.
\end{euclidtheorem}
\begin{proof}
This is, of course, where you include the construction and proof.  There should be diagrams!
\end{proof}

% PROBLEM #2 - RESULT OF SOMEONE ELSE
\item  The previous result that Sarah asked me to prove was that the diagonals of a kite are perpendicular.
\setcounter{section}{2}
\setcounter{theorem}{4}
% Change the introduction and counters appropriately to reflect the problem I asked you to share.
\begin{theorem}
If the diagonals of a kite meet, then they meet at a right angle.
\end{theorem}
\begin{proof}
Put the proof here.  It should be your own proof, not simply a copy of the work of the original presenter or author of the paper.  It does not need to be exactly the same as the proof you gave during your oral exam.
\end{proof}

% PROBLEM #3 - PARALLELOGRAMS
\item  I determined that the conjecture was true.
% If you determined it was false, change that.
% The * here means that the theorem will not have a number.
\begin{theorem*}
The diagonals of a parallelogram bisect one another.
\end{theorem*}
% If appropriate, change the statement of the theorem to reflect the conjecture is false.
\begin{proof}
Put the proof of the theorem here.
\end{proof}

% PROBLEM #4 - 3.7
\item  I determined the the conjecture was false.
% If you determined it was true, change that.
\setcounter{section}{3}
\setcounter{theorem}{6}
% You do not need to adjust these counters.
\begin{theorem}
% If you proved the conjecture false, the theorem may be something like
Let $ABCD$ be a quadrilateral.  The midpoints of the four sides need not form the vertices of a parallelogram.
\end{theorem}
\begin{proof}
Include the proof of the resulting theorem here.
%If you have determined the conjecture is false, your proof may be constructive.  i.e.  The following is the construction of a non-convex quadrilateral where the midlines $EF$ and $GH$ are not parallel.
\end{proof}


% PROBLEM #5 - CONTRADICTION
\item  Include here a description of the process you used to determine which result you would prove via contradiction.
\setcounter{section}{1}
\setcounter{theorem}{4}
%change these counters so that the numbering is correct.
\begin{euclidtheorem}
% Choose whichever theorem environment reflects the result you chose to prove by contradiction.

%\begin{theorem}

In triangles $ABC$ and $DEF$, if sides $AB$ and $DE$ are congruent, sides $BC$ and $EF$ are congruent, and angles $ABC$ and $DEF$ are congruent, then the triangles and corresponding parts are congruent.

%\end{theorem}
\end{euclidtheorem}
% You may restate a result of Euclid in more modern language, particularly if is helps to make your proof by contradiction make more sense.

\begin{proof}
Be extra clear in your proof set-up and in the contradiction that you reach.
\end{proof}

\end{enumerate}

\subsection*{Paper Folding}

For this portion, you should be writing for an audience that is unfamiliar with this question.  Explain the folding required 
% I've copied my explanation below in case that is helpful, but definitely don't just use mine.
and any special cases you noticed where the folding directions should be different.  Some results are difficult to state formally.  Include those types of things that you noticed in this narrative portion.

In a second part, include formal statements of any conjectures you have.  You should aim to move as many statements from the informal section above to this part as you can.

In the third section, include the proof of one of your conjectures.  Again, the more you have, the better, but definitely prove something.

In a final section, include any related ideas you could explore.  You don't have to actually do any of this nor make any formal or informal conjectures.  These statements could very likely start with the phrase, ``I wonder what would happen if \dots''



Start with a square piece of paper.  Make a single fold in the paper.  Any fold you wish.
\begin{center}\includegraphics[width = .35\textwidth]{square} \includegraphics[width = .35\textwidth]{motherline}\end{center}
\noindent This splits the four edges of the square into six segments.  Fold each of these pieces to meet the main blue fold:
\begin{center}\includegraphics[width = .35\textwidth]{firstfold}\includegraphics[width = .35\textwidth]{secondfold}\includegraphics[width = .35\textwidth]{allfolds}\end{center}


\end{document}